% ============================================================
% APPENDIX G -- Data sources and their verification
% ============================================================

\label{app:datachecks}

\begin{ataglance}[How to read this appendix]
This appendix records how the NHATS variables and the health-utility tariff
used in the analysis were checked against their published sources, what the
checks changed, and what they could not settle. \S\ref{sec:G_sources} lists
the documents consulted. \S\ref{sec:G_names} reports the check of variable
names and files, and \S\ref{sec:G_codes} the check of response codes.
\S\ref{sec:G_changes} sets out the corrections that followed.
\S\ref{sec:G_tariff} reports the check of the EQ-5D tariff.
\S\ref{sec:G_limits} states the shortcomings that remain and how each bears on
the results. \S\ref{sec:G_enforce} describes how the checks are enforced when
the analysis is run.
\end{ataglance}

\section{Purpose and Sources}
\label{sec:G_sources}

The analytic file is built by a script from the NHATS public-use files, and
every variable name and response code it uses is declared in one layout file
(Section~\ref{sec:survival}). The script can confirm that a declared variable
exists in the files supplied to it, and that each observed value is one of the
declared codes. It cannot confirm that a declared code means what the layout
assumes it means. That check has to be made against the documents NHATS
publishes, and it was made, item by item, on 27 October 2026, before any
analysis of NHATS data.

Table~\ref{tab:G_sources} lists the documents consulted and what each was used
to check. The Round~5 crosswalk lists every public-use variable together with
the file that holds it and the sample groups to which the question was put
\citep{nhats_r5_crosswalk_2024}. The Round~5 instrument gives the wording and
answer codes of each question and the routing that decides who is asked
\citep{nhats_r5_instrument_2024}. The User Guide defines the variables NHATS
derives from several questions, and records known errors in the released files
\citep{freedman_nhats_userguide_2026}. Three further NHATS documents settled
questions the first three left open: the Round~5 data collection procedures,
which list every final case status code \citep{nhats_dcp_r5_2015}; the
Round~10 weighting paper, which confirms that the same codes are used in
later rounds \citep{dematteis_2021_r10weights}; and the income imputation
paper, which gives the name of the auxiliary income file
\citep{hu_freedman_2024_income}. The Round~6 crosswalk identified the public
variables that date the outcomes \citep{nhats_r6_crosswalk_2024}. The tariff was checked against the
discussion paper in which the UK time trade-off tariff was first published
\citep{dolan_1995_che}. The tariff is more often cited from the later journal
article \citep{dolan_1997}, which was not available for this check.

\begin{table}[htbp]
\centering
\caption{Documents against which the data layout was checked.}
\label{tab:G_sources}
\footnotesize
\renewcommand{\arraystretch}{1.2}
\begin{tabular}{p{5.2cm} p{2.6cm} p{6.2cm}}
\toprule
\textbf{Document} & \textbf{Version} & \textbf{Used to check} \\
\midrule
NHATS Round 5 crosswalk between instruments and codebook (spreadsheet) & 15 February 2024 & Variable names; the file that holds each variable; which sample groups were asked \\
NHATS Round 5 data collection instrument & February 2024 & Question wording, answer codes and routing \\
NHATS User Guide, Rounds 1 to 14 & August 2026 & Codes of derived variables; residence status; data notices \\
NHATS Data Collection Procedures, Round 5 & 2015 & Final case status codes (Table II.F-1) \\
NHATS Technical Paper 28, Round 10 weights & July 2021 & Status codes in a later round (Table 1) \\
NHATS Technical Paper 43, income imputations & May 2025 & Name and content of the auxiliary income file \\
NHATS Round 6 crosswalk between instruments and codebook & 15 February 2024 & Public variables that date interviews and deaths \\
Dolan, Gudex, Kind and Williams (1995), Discussion Paper 138 & 1995 & All coefficients of the UK EQ-5D time trade-off tariff \\
\bottomrule
\end{tabular}
\end{table}

\section{Variable Names and Files}
\label{sec:G_names}

Every variable named in the layout was looked up in the Round~5 crosswalk.
Three names were wrong. All three are items of the physical capacity section
used in the frailty comparator, and had been written from the NHATS naming
pattern rather than read from the crosswalk. Their correct names were found in
the same section of the crosswalk and the layout was corrected
(Table~\ref{tab:G_changes}). After the corrections described below, all 89
Round~5 variables in the layout appear in the Round~5 crosswalk, and the two
Round~1 variables it draws on appear in the Round~1 crosswalk.
The design variables (analytic weight, variance stratum and variance cluster)
are on the sample-person file, as are the derived race and ethnicity variable
and the marital status used for income equivalisation. Vital status is on the
tracker file. Total income is held in three forms: the reported amount and
five hot-deck imputations on the sample-person file, and twenty
interval-regression imputations in a separate income imputation file
\citep{freedman_nhats_userguide_2026}.

The crosswalk also shows which sample groups were asked each question, and
this exposed a gap that a check of names alone would not find. The education
question, and the question from which gender is derived, are put only to people
at their first interview. Continuing members of the 2011 cohort answered them
in Round~1, so their Round~5 values are missing by design. Without a further
step, education would have been missing for every continuing member of the
analytic cohort. The build now fills these
two items from the Round~1 sample-person file for continuing members. The
derived race and ethnicity variable needs no such step, since NHATS carries the
Round~1 value forward into the Round~5 file \citep{freedman_nhats_userguide_2026}.

\section{Response Codes}
\label{sec:G_codes}

Table~\ref{tab:G_codes} summarises the codes read from the instrument and the
User Guide for each indicator, and how each code is used. For most items the
codes are the NHATS convention for a yes-or-no question (1 yes, 2 no) and
matched the layout. Three groups of items needed more care.

The twelve activity items behind the function indicator are not answers to a
single question. NHATS derives each from several questions and codes the result
as follows: 1, did not do the activity by oneself in the last month; 2, did it
by oneself with no difficulty; 3, did it by oneself with difficulty; 4 and 5,
as 2 and 3 when whether it was done alone is unknown; 6 and 7, difficulty
unknown; 8, not done in the last month; and, for medications only, 9, takes no
medications \citep{freedman_nhats_userguide_2026}. The layout had read these
as a scale of increasing difficulty, with 1 meaning no difficulty. Under that
reading code~2, which means no difficulty, was counted as a limitation. The
indicator now counts codes 1, 3 and 5 as a limitation, codes 2, 4 and 9 as
none, and treats codes 6, 7 and 8 as missing.

For people continuing from earlier rounds, the Round~5 disease questions ask
only about conditions diagnosed since the last interview, and a condition
already reported is coded 7 \citep{nhats_r5_instrument_2024}. A condition is
therefore counted as present when coded 1 or 7, which is how the layout
already read it; the check confirmed the reading rather than changing it.

Two home environment questions are skipped by design for some people. The
number of floors is not asked in multi-unit buildings, and whether bedroom,
kitchen and bathroom share a floor is not asked of people in one-floor homes
\citep{nhats_r5_instrument_2024}. The skipped answer carries the NHATS code for
``inapplicable'', which the build would otherwise treat as missing. For a
one-floor home the skipped answer is known: the rooms share a floor. The build
reads both skips as the absence of a barrier. The floors question itself has
four answers, the last being four or more, not five as the layout had assumed.

\begin{table}[htbp]
\centering
\caption{Response codes used for each indicator, as read from the Round~5 instrument and the User Guide.}
\label{tab:G_codes}
\footnotesize
\renewcommand{\arraystretch}{1.15}
\begin{tabular}{p{3.0cm} p{5.6cm} p{5.4cm}}
\toprule
\textbf{Indicator} & \textbf{Codes} & \textbf{Use} \\
\midrule
Chronic conditions & 1 yes; 2 no; 7 previously reported & Present if 1 or 7 \\
Function & 1 to 9, derived (see text) & Limitation if 1, 3 or 5; missing if 6 to 8 \\
Cognition & Recall 0 to 10; clock 0 to 5 & Summed; reserved codes missing \\
Depression and anxiety & 1 not at all to 4 nearly every day & Reversed and summed \\
Participation & 1 yes; 2 no (last month) & Count of yes \\
Social network & 0 to 5 people & Value; not asked of proxies \\
Communication & 1 yes; 2 no & Count of yes \\
Income & Dollars; marital status 1 to 6 & Minus mean log equivalised income \\
Education & 1 no schooling to 9 master's or doctoral & Reversed \\
Home access & Stairs, ramp, same floor 1 yes, 2 no; floors 1 to 4 & Count of barriers; skips read as none \\
Care access & 1 yes; 2 no & Count of no \\
\bottomrule
\end{tabular}
\end{table}

\section{Corrections Made}
\label{sec:G_changes}

Table~\ref{tab:G_changes} lists every change made to the layout or the build
as a result of the checks, with the source that required it. Two of them would
have changed the results had they gone unnoticed. The reversed activity codes
would have counted people who manage an activity alone without difficulty as
limited in it, and so inverted part of the function indicator. The missing
education values would have left one of the four systemic indicators to be
imputed for every continuing member of the cohort. A third
concerns the sample itself. The User Guide records that people who died before
the Round~5 interview were given a positive 2015 analytic weight in error
\citep{freedman_nhats_userguide_2026}. They have no Round~5 interview and
cannot contribute baseline indicators; the build now restricts the analytic
sample to people with a Round~5 sample-person interview, which removes them
along with facility-only cases.

\begin{table}[htbp]
\centering
\caption{Changes to the data layout and build that followed from the checks.}
\label{tab:G_changes}
\footnotesize
\renewcommand{\arraystretch}{1.15}
\begin{tabular}{p{4.2cm} p{6.4cm} p{3.4cm}}
\toprule
\textbf{Item} & \textbf{Change} & \textbf{Source} \\
\midrule
Function indicator and two comparators & Activity codes read as in the User Guide; code 2 is no limitation & User Guide, derived variables \\
Education & Filled from the Round~1 file for continuing members & Crosswalk; User Guide \\
Gender & Filled from the Round~1 file where missing & Crosswalk \\
Analytic sample & Restricted to people with a Round~5 interview; decedents excluded & User Guide, residence status and data notices \\
Frailty comparator & Three names corrected; reach overhead and bend over both used & Crosswalk \\
Floors & Four codes, not five & Instrument \\
Home access & Skipped questions read as no barrier & Instrument routing \\
Race and ethnicity & Code 5 (more than one race, no primary) as other; code 6 missing & User Guide \\
Income & Twenty imputations used when present, else five; couple income divided by $\sqrt{2}$ & Crosswalk; User Guide \\
Income file & Named as NHATS names it; a renamed file for the round is also found & Technical Paper 43 \\
Deaths & Status 86 (deceased, no proxy) counted with 62 & Data Collection Procedures \\
Follow-up & Extended only by a status showing the person alive & Data Collection Procedures \\
Imputation model & Proxy respondent, cognition not asked, reported dementia and the AD8 proxy screen added & Instrument; User Guide \\
Event times & Dated in months from interview, case-status and death-interval variables, not whole rounds & Round 6 crosswalk; User Guide \\
\bottomrule
\end{tabular}
\end{table}

The changes to the mortality outcome need a word of explanation. NHATS
assigns each case a final status at every round. Two codes record a death:
62, when a proxy completed the last-month-of-life interview, and 86, when
the person had died and no proxy was available \citep{nhats_dcp_r5_2015}.
The layout had listed only 62, so deaths with no proxy would have been
treated as losses to follow-up. Both now count. The same table showed a
second fault. The build had treated any recorded status as evidence that the
person was alive in that round, so a case coded ``unable to locate'' extended
follow-up to a round in which the person might already have died. Follow-up
now extends only through rounds with a status that shows the person was
reached or known to be alive: a completed interview, a refusal or inability
to take part, absence during fieldwork, a move outside the area or out of the
country, or a refusal by the facility. A proxy's refusal is not taken as a
sign of life, because the proxy may be refusing an interview about a death.

The timing of events changed as well. The first version placed every event at
the midpoint of a whole round, so that all deaths fell at one of five times.
The public files carry more than that. The tracker gives the month and year of
each sample-person interview, with a revised month for Rounds 3 to 8, and the
month each case status was set. For a death with a last-month-of-life
interview, the follow-up file gives the number of months between the last
interview and the death \citep{freedman_nhats_userguide_2026}. Times are now
measured in months from the Round~5 interview. A death with a proxy interview
is dated by the reported interval; a death without one, by the midpoint
between the last month the person was known alive and the month the death was
recorded. A first hospital stay is placed midway between the interview that
reports it and the previous one. The exact date of death is a restricted
variable and was not used.

The income change needs a word of explanation as well. The previous layout used the
reported total, which is missing for a large share of respondents and, for
those married or living with a partner, is the income of the couple. The build
now averages the logarithm of income over the NHATS imputations and divides
couple income by the square root of two, the conventional square-root
equivalence scale, so that a single person and a couple with the same income
per head are scored alike.

\section{The Health-Utility Tariff}
\label{sec:G_tariff}

The utility comparator values crosswalked EQ-5D levels with the UK time
trade-off tariff. Each of its eleven coefficients was compared with Table~1 of
\citet{dolan_1995_che}, and all eleven agree (Table~\ref{tab:G_tariff}). Two
further checks were made by computation. The paper's own worked example, state
11223, gives $1 - 0.081 - 0.036 - 0.123 - 0.236 - 0.269 = 0.255$, as printed.
The worst state, 33333, gives $-0.594$; since every coefficient is positive,
no state can score lower.

\begin{table}[htbp]
\centering
\caption[UK EQ-5D time trade-off coefficients used in the utility comparator]{UK EQ-5D time trade-off coefficients used in the utility comparator, as checked against \citet{dolan_1995_che}, Table~1.}
\label{tab:G_tariff}
\footnotesize
\renewcommand{\arraystretch}{1.15}
\begin{tabular}{p{5.0cm} p{2.0cm} p{2.0cm}}
\toprule
\textbf{Term} & \textbf{Level 2} & \textbf{Level 3} \\
\midrule
Mobility & 0.069 & 0.314 \\
Self-care & 0.104 & 0.214 \\
Usual activities & 0.036 & 0.094 \\
Pain or discomfort & 0.123 & 0.386 \\
Anxiety or depression & 0.071 & 0.236 \\
\midrule
Constant (any level above 1) & \multicolumn{2}{l}{0.081} \\
N3 (any level 3) & \multicolumn{2}{l}{0.269} \\
\bottomrule
\end{tabular}
\end{table}

The tariff being right does not make the comparator a validated utility. NHATS
does not administer the EQ-5D, and the levels are assigned from NHATS items by
a crosswalk written for this thesis. The comparator is labelled a utility
proxy throughout for that reason.

\section{Shortcomings That Remain}
\label{sec:G_limits}

The checks settled most questions about names and codes. The following were
not settled, or cannot be settled from documents, and the reader should weigh
the results with them in mind.

\paragraph{Timing of deaths.} A death with a proxy interview is dated to the
month, from the interval the proxy reports. A death without one, which is the
smaller group, is known only to lie between the last month the person was
seen alive and the month the death was recorded, and is placed at the
midpoint. A first hospital stay is likewise known only to lie between two
interviews. Errors of this kind are of the order of a few months and shift
event times in both directions; they matter little for five-year risk but
more for risk at one year.

\paragraph{One name differs between documents.} The derived medication variable
is listed in the crosswalk as \texttt{mc5dmedssfdf} and in the User Guide as
\texttt{mc\#dmedsfdf}. The layout uses the crosswalk spelling. The build stops
if the variable is absent, so a mismatch cannot pass unnoticed.

\paragraph{Proxy respondents.} Social network size is not asked when a proxy
answers for the sample person, and word recall and clock drawing are not
administered when the proxy says the person cannot be asked. These values are
missing for reasons tied to the person's health. Four baseline variables are
therefore added to the imputation model: whether a proxy answered, whether
the proxy said the cognitive tests could not be given, whether dementia has
been reported, and the score on the AD8, the eight-item dementia screen NHATS
puts to proxy respondents, scored as NHATS recommends
\citep{freedman_nhats_userguide_2026}. The flags are never outcomes, so the measurement stage
remains blind to them. They help where people in the same position were
measured: proxy cases whose cognition was tested, for instance, inform the
imputation of other proxy cases. The AD8 matters most here. A flag that is set
only for people with no cognition score cannot teach the model anything about
their cognition, since no one in that group has a score to learn from. The AD8
can: proxies answer it both for people who were tested and for people who were
not, so the model learns how test scores fall as the AD8 rises among the
tested, and carries that relation to the untested. That relation is still an
extrapolation, because the untested sit at the severe end of the AD8, beyond
most of the people from whom it was learned. A sensitivity analysis therefore
sets their cognition to the lowest possible score instead. Chapter~\ref{ch:results1}
reports how far the loadings and the biomedical domain scores move between
the two rules. The truth lies between them: the imputation is likely to
overstate the cognition of these people, and the lowest score to understate
it for some. NHATS classifies dementia in such cases from a proxy screening
instrument; that classification was not used as an indicator, because it
would bring a diagnostic judgement into an indicator meant to measure
performance.

\paragraph{Income imputation.} Averaging the logarithm of income over the NHATS
imputations gives one value per person and discards the spread between
imputations. Standard errors that involve the income indicator are therefore
somewhat too small. The square-root scale for couples is a convention, not an
estimate.

\paragraph{Skipped home items.} Reading the skipped floors question as one
level is exact for one-floor homes but an assumption for multi-unit buildings,
where a dwelling may span more than one level. The effect is to understate
barriers for a minority of people in multi-unit buildings.

\paragraph{First-interview items.} Education and gender for continuing members
are taken from Round~1, four years before the index assessment. Education
completed by people aged 65 and over is unlikely to change in that interval.
A continuing member whose Round~1 answer is itself missing stays missing.

\paragraph{Documents and release.} The User Guide consulted is the edition
covering Rounds 1 to 14. Where NHATS has reissued a Round~5 file since its
first release, the names and codes may differ from those documented at the
time. The pre-analysis check described next guards against that, but only for
names and codes, not for changes of meaning.

\section{How the Checks Are Enforced}
\label{sec:G_enforce}

The checks above are recorded in two places that the analysis cannot bypass.
The layout file cites, beside each variable, the document from which its name
and codes were read. Before the analytic file is built, the pipeline confirms
that every declared variable is present in the files supplied, reports every
value outside the declared codes and every tracker status value, and records
which income source it used, and stops if
any required variable is absent. The test suite of the pipeline contains a
test for each correction in Table~\ref{tab:G_changes}: that code~2 of an
activity item counts as no limitation, that a skipped home question counts as
no barrier, that education and gender are filled from Round~1, that decedents
leave the sample, that income is taken from the twenty imputations and
equivalised and that a renamed income file is still found, that a death with
no proxy counts as a death, that an untraced round does not extend follow-up,
and that the proxy flags and the AD8 score reach the imputation model, that the
AD8 is scored as NHATS recommends, and that deaths and hospital stays are
dated from the NHATS months. These tests run on every change to the code, so a later edit
that reintroduced one of the errors would be caught before it reached the
results.
